\chapter{A Normalized Internal Record Process}
\label{app:records}

This appendix constructs a finite-history law from an internal interaction rule. The rule is supplied; normalization does not select it. Let $\mathsf S$ and $\mathsf O$ be standard Borel spaces. A state $x$ includes an accessible record register. Write $\rho:\mathsf S\to\mathsf O$ for its readout.

\constructionheading{Event and registration}{def:registered-event}
Choose a normalized kernel $R:\mathsf S\rightsquigarrow\mathsf O$ and a measurable update $A:\mathsf S\times\mathsf O\to\mathsf S$ satisfying
\begin{equation}
\rho(A(x,o))=o,
\label{eq:registration-condition}
\end{equation}
for $R(x,\cdot)$-almost every $o$, for every $x$ in a Borel admitted set $\mathsf S_{\mathrm{adm}}\subseteq\mathsf S$. Define
\begin{equation}
J(x,dx',do)=R(x,do)\delta_{A(x,o)}(dx').
\label{eq:registered-transition}
\end{equation}
If necessary, the state can be explicitly augmented by a register so that \cref{eq:registration-condition} holds. The admitted set is assumed to contain $\mu_0$-almost every initial state and to be invariant, $J(x,\mathsf S_{\mathrm{adm}}\times\mathsf O)=1$ for $x\in\mathsf S_{\mathrm{adm}}$, so that the registration condition holds at every visited state. The map $x\mapsto J(x,B)$ is measurable for every measurable $B\subseteq\mathsf S\times\mathsf O$, because $x\mapsto\int f(x,o)R(x,do)$ is measurable for every bounded jointly measurable $f$, by a monotone-class argument from products. When the event kernel is generated by an interaction transition $\Gamma$ and the readout, $R$ is the record kernel of \cref{prop:record-law}, and the state marginal of $J$ reproduces it: $R(x,B)=\int\mathbf 1_B(\rho(A(x,o)))R(x,do)$.

\propositionheading{Consistent finite histories}{prop:registered-history}
For an initial law $\mu_0$ on $\mathsf S$, the iterated kernel product
\begin{equation}
\mathbb P_n(dx_0,do_1,dx_1,\ldots,do_n,dx_n)
=\mu_0(dx_0)\prod_{t=1}^nJ(x_{t-1},dx_t,do_t),
\label{eq:registered-history}
\end{equation}
is a probability law. Marginalizing its last pair gives $\mathbb P_{n-1}$, and every record agrees with the readout of the resulting state almost surely.
\begin{proof}
The product denotes successive kernel integration, not independent factors. Integrating the last pair gives one because $J(x,\mathsf S\times\mathsf O)=1$. Repeating this argument gives total mass one and the stated consistency. Each transition is supported on $x'=A(x,o)$ and hence on $\rho(x')=o$. For a finite history, the intersection of these probability-one events still has probability one.
\end{proof}

By the Ionescu-Tulcea theorem \citep{Kallenberg2021}, the consistent family $(\mathbb P_n)_n$ extends to a unique law on infinite histories, on which persistence statements about later readouts can be posed. When each record is binary, the records of a finite history carry at most $n\log2$ nats about any variable, since $I(X;O_{1:n})\leq H(O_{1:n})\leq n\log2$.

Nothing in this construction requires a variable external to the agents. It does require a declared accessible register and an update. A stochastic $A$ is handled by replacing its Dirac kernel with a normalized kernel supported on the registration condition. A deterministic transition is recovered when $R$ is a Dirac kernel as well.

\section{Prediction followed by assimilation}

One possible internal state contains a model law $s$ and a record register. Fix a normalized density $k(o\mid m)$ with respect to a reference measure on $\mathsf O$. Let
\begin{equation}
R(s,do)=\left[\int k(o\mid m)s(dm)\right]do,
\qquad
s^o(dm)=\frac{k(o\mid m)s(dm)}{\int k(o\mid m')s(dm')}.
\label{eq:predict-and-assimilate}
\end{equation}
The update replaces $s$ by $s^o$ and stores $o$. The denominator is positive and finite for predictive-almost every record. Defining an arbitrary probability law on the remaining null set completes the measurable rule under the usual conditional-law hypotheses.

\propositionheading{The Bayesian history identity}{prop:bayesian-history}
For a fixed normalized observation density $k$ and initial law $s_0$, sequential prediction and Bayesian assimilation give the record density
\begin{equation}
p(o_1,\ldots,o_n)
=\int\prod_{t=1}^n k(o_t\mid m)s_0(dm).
\label{eq:bayesian-history}
\end{equation}
\begin{proof}
After $t-1$ records, the posterior is proportional to $s_0(dm)\prod_{r<t}k(o_r\mid m)$. The next predictive density is the integral with the product extended to $t$, divided by the integral up to $t-1$. Multiplying these conditional densities cancels every intermediate normalizer and leaves \cref{eq:bayesian-history}.
\end{proof}

Let the presentations $m$ be the laws on a finite record alphabet, let $k(o\mid m)=m(o)$, and let $s_0$ be a Dirichlet law with parameter $\alpha$ and total concentration $A$. The predictive law after $n$ records is $\bar m_n=(\alpha+c_n)/(A+n)$, with $c_n$ the record counts. Then $\bar m_n-\bar m_{n-1}=(e_{o_n}-\bar m_{n-1})/(A+n)$, so the total-variation step is at most $1/(A+n)$: a slow model sector produced by accumulated records rather than by an inserted rate.

This identity applies to the stated fixed observation model. If the likelihood itself is changed using the current state, the product of the actual predictive kernels still defines a normalized process, but it need not equal the fixed-latent expression. Likewise, replacing Bayes' rule by a finite natural-gradient step defines a different assimilation mechanism unless the updates coincide. The normalization of a process and its interpretation as Bayesian conditioning are separate statements. Their difference has an exact measure.

\propositionheading{Departure from Bayesian assimilation}{prop:variational-history}
Let $\mathbb P_n$ be the record law of $n$ steps of \cref{eq:predict-and-assimilate} and let $\tilde{\mathbb P}_n$ be the record law generated from the same $s_0$ and $k$ by any measurable assimilation rule, with model laws $s_{t-1}$ and $\tilde s_{t-1}$ after the same history $o_{1:t-1}$. Then
\begin{equation}
\KL(\mathbb P_n\Vert\tilde{\mathbb P}_n)
=\sum_{t=1}^n\E_{\mathbb P}\KL(s_{t-1}k\Vert\tilde s_{t-1}k)
\leq\sum_{t=1}^n\E_{\mathbb P}\KL(s_{t-1}\Vert\tilde s_{t-1}),
\label{eq:variational-history}
\end{equation}
where $s_{t-1}k$ denotes the predictive record law.
\begin{proof}
Both processes predict the next record by $s_{t-1}k$ and $\tilde s_{t-1}k$ respectively, so the chain rule for relative entropy over the successive records gives the equality, and the data-processing inequality for the kernel $k$ gives each bound.
\end{proof}

The identity applies to record laws. The joint laws of states and records of the two processes are generally mutually singular, because their state updates are different deterministic functions of the same history.
